\section{Gambaran Umum Program}
Implementasi Milestone 4 melanjutkan pipeline Milestone 1, Milestone 2, dan Milestone 3. Alur utama berada pada \texttt{src/main.cpp}. Program membaca file input, lalu menentukan bentuk input tersebut. Jika input adalah parse tree, program menggunakan \texttt{ParseTreeReader}. Jika input adalah token stream, program menggunakan \texttt{TokenStreamReader} lalu menjalankan parser. Jika input adalah source code Arion, program menjalankan lexer dan parser.

Setelah parse tree tersedia, program menjalankan \texttt{ASTBuilder} untuk membentuk AST. AST tersebut kemudian dianalisis oleh \texttt{SemanticAnalyzer}. Jika tidak ada error, program menjalankan \texttt{IntermediateCodeGenerator} untuk mengubah decorated AST menjadi daftar instruksi intermediate code. Terakhir, \texttt{Interpreter} mengeksekusi intermediate code tersebut dan menghasilkan output final.

\begin{figure}[H]
    \centering
    \begin{tikzpicture}[
        scale=0.88,
        transform shape,
        node distance=8mm and 18mm,
        every node/.style={draw, rounded corners, align=center, minimum width=3.4cm, minimum height=7mm, font=\small},
        arrow/.style={-{Stealth[length=2.5mm]}, thick}
        ]
        \node (input) {Input file};
        \node (lex) [below=of input] {Lexer / TokenStreamReader / ParseTreeReader};
        \node (parser) [below=of lex] {Parser};
        \node (parsetree) [below=of parser] {Parse Tree};
        \node (astbuilder) [below=of parsetree] {ASTBuilder};
        \node (ast) [below=of astbuilder] {AST};
        \node (sem) [below=of ast] {SemanticAnalyzer};
        \node (codegen) [below=of sem] {IntermediateCodeGenerator};
        \node (interpreter) [below=of codegen] {Interpreter};
        \node (out) [below=of interpreter] {Output Final};

        \draw[arrow] (input) -- (lex);
        \draw[arrow] (lex) -- (parser);
        \draw[arrow] (parser) -- (parsetree);
        \draw[arrow] (parsetree) -- (astbuilder);
        \draw[arrow] (astbuilder) -- (ast);
        \draw[arrow] (ast) -- (sem);
        \draw[arrow] (sem) -- (codegen);
        \draw[arrow] (codegen) -- (interpreter);
        \draw[arrow] (interpreter) -- (out);
    \end{tikzpicture}
    \caption{Pipeline Program Milestone 4}
    \label{fig:pipeline-program-m4}
\end{figure}

\section{Struktur Direktori Program}
Bagian Milestone 4 terutama berada pada direktori \texttt{src/codegen/} dan \texttt{src/interpreter/}. Berkas utama yang digunakan adalah sebagai berikut.

\begin{itemize}
    \item \texttt{intermediate\_code.h} dan \texttt{intermediate\_code.cpp}: definisi instruction set, intermediate code generator, dan konversi decorated AST menjadi daftar instruksi.
    \item \texttt{interpreter.h} dan \texttt{interpreter.cpp}: implementasi virtual machine dengan siklus fetch-decode-execute, stack management, dan eksekusi instruksi.
    \item \texttt{runtime\_value.h} dan \texttt{runtime\_value.cpp}: representasi nilai runtime menggunakan \texttt{std::variant} untuk mendukung \texttt{int}, \texttt{double}, \texttt{char}, \texttt{bool}, dan \texttt{std::string}.
\end{itemize}

\section{Implementasi Intermediate Code Generator}
Intermediate code generator diimplementasikan pada class \texttt{CodeGen::IntermediateCodeGenerator}. Fungsi utama \texttt{generate} menerima \texttt{ProgramNode} dari decorated AST dan \texttt{SymbolTable} hasil semantic analysis, kemudian menghasilkan \texttt{Result} yang berisi daftar instruksi dan error (jika ada).

Generator bekerja secara top-down, mengunjungi setiap node AST dan mengeluarkan instruksi yang relevan. Detail implementasi meliputi beberapa aspek berikut.

\subsection{Inisialisasi Memori}
Fungsi \texttt{initialFrameSize} menghitung ukuran frame awal berdasarkan informasi \texttt{btab} dari program. Frame size meliputi \texttt{kFrameHeaderSize} (3 slot untuk static link, dynamic link, dan return address) ditambah jumlah parameter \texttt{(psze)} dan variabel lokal \texttt{(vsze)}. Instruksi \texttt{INT} pertama yang dihasilkan adalah \texttt{INT 0 frameSize}.

\subsection{Translasi Statement}
Setiap jenis statement memiliki fungsi \texttt{visit} sendiri.

\begin{itemize}
    \item \textbf{Assignment:} Generator mengevaluasi ekspresi value terlebih dahulu, lalu memanggil \texttt{emitStore} untuk target. Target dapat berupa variabel sederhana atau array access.
    \item \textbf{If:} Kondisi dievaluasi, diikuti \texttt{JPC} ke label else/end. Blok then diterjemahkan, lalu \texttt{JMP} ke end. Jika ada else, blok else diterjemahkan, lalu label end dipasang.
    \item \textbf{While:} Label loop start dipasang, kondisi dievaluasi, \texttt{JPC} ke end, body diterjemahkan, \texttt{JMP} kembali ke start, lalu label end dipasang.
    \item \textbf{Repeat:} Label loop start dipasang, body diterjemahkan, kondisi dievaluasi, \texttt{JPC} kembali ke start.
    \item \textbf{For:} Variabel loop diinisialisasi dengan start expression. Label loop start dipasang, variabel loop dan end expression dievaluasi, dibandingkan dengan \texttt{LEQ} atau \texttt{GEQ}, \texttt{JPC} ke end, body diterjemahkan, variabel loop diincrement/decrement, \texttt{JMP} kembali ke start, lalu label end dipasang.
    \item \textbf{Case:} Untuk setiap branch, selector dan label dievaluasi, dibandingkan dengan \texttt{EQL}, \texttt{JPC} ke label berikutnya jika tidak cocok. Jika cocok, eksekusi jatuh ke body branch. Setelah body, \texttt{JMP} ke end case. Semua \texttt{JMP} ke end di-backpatch setelah seluruh branch diproses.
\end{itemize}

\subsection{Translasi Ekspresi}
Ekspresi diterjemahkan secara rekursif. Literal integer, real, char, string, dan Boolean menghasilkan \texttt{LIT}. Variabel menghasilkan \texttt{LOD}. Array access dengan satu atau dua dimensi menghasilkan \texttt{ALOD} dengan operand yang menyimpan metadata dimensi. Operator biner menghasilkan evaluasi kedua operand diikuti \texttt{OPR}. Unary minus menghasilkan evaluasi operand diikuti \texttt{OPR NEG}.

\subsection{Translasi Procedure Call}
Saat memanggil prosedur, generator mengevaluasi argumen (push ke stack), lalu mengeluarkan \texttt{CAL} dengan target alamat prosedur. Jika alamat belum diketahui saat pemanggilan, instruksi \texttt{CAL} menggunakan placeholder yang di-backpatch setelah seluruh deklarasi prosedur diproses. Setelah \texttt{CAL} selesai, generator mengeluarkan \texttt{INT} dengan nilai negatif untuk membersihkan argumen dari stack caller.

\subsection{Backpatching}
Beberapa instruksi seperti \texttt{JMP}, \texttt{JPC}, dan \texttt{CAL} memerlukan backpatching karena alamat tujuan belum diketahui saat instruksi pertama kali dihasilkan. Generator menyimpan indeks instruksi yang perlu di-backpatch dalam variabel sementara, lalu mengisi operandnya setelah alamat tujuan diketahui.

\section{Implementasi Interpreter}
Interpreter diimplementasikan pada class \texttt{Interpreter::Interpreter}. Fungsi utama \texttt{run} menerima daftar instruksi, menginisialisasi stack dan register, lalu mengeksekusi instruksi satu per satu dalam siklus fetch-decode-execute.

\subsection{Siklus Eksekusi}
Interpreter menggunakan variabel \texttt{ip} (instruction pointer), \texttt{bp} (base pointer), dan \texttt{sp} (stack pointer). Siklus berjalan selama \texttt{ip} masih dalam rentang valid dan tidak ada error.

\begin{breakablealgorithm}[Siklus Fetch-Decode-Execute]
    \label{alg:fde-cycle}
    \begin{algorithmic}[1]
        \Procedure{Run}{program}
        \State \texttt{reset()}
        \State \texttt{validateJumpTargets(program)}
        \While{\texttt{not halted} \textbf{and} \texttt{ip} valid}
            \State \texttt{inst = program[ip]}
            \State \texttt{advance = true}
            \If{\textbf{not} \Call{executeInstruction}{inst, program, advance}}
                \State \Return \texttt{error}
            \EndIf
            \If{\texttt{advance}}
                \State \texttt{ip += 1}
            \EndIf
        \EndWhile
        \State \Return \texttt{result}
        \EndProcedure
    \end{algorithmic}
\end{breakablealgorithm}

\subsection{Manajemen Stack}
Stack diimplementasikan menggunakan \texttt{std::vector<RuntimeValue>}. Fungsi \texttt{push} menambahkan nilai ke stack dengan pengecekan batas maksimal \texttt{kMaxStackSlots}. Fungsi \texttt{pop} mengambil nilai dari stack dengan pengecekan underflow. Fungsi \texttt{ensureStackSize} mengubah ukuran vector sesuai kebutuhan.

Untuk instruksi \texttt{INT} dengan nilai negatif, interpreter menyusutkan stack dengan menghapus elemen dari belakang, asalkan tidak menyusut di bawah frame header saat ini.

\subsection{Frame Management}
Instruksi \texttt{CAL} membuat frame baru dengan menyimpan static link, dynamic link, dan return address. Static link dihitung berdasarkan lexical level delta. Dynamic link adalah \texttt{bp} saat ini. Return address adalah \texttt{ip + 1}. Instruksi \texttt{RET} memulihkan \texttt{ip} dari return address, \texttt{bp} dari dynamic link, dan menghapus frame saat ini.

\subsection{Validasi Keamanan}
Interpreter memvalidasi target jump sebelum eksekusi dimulai untuk memastikan semua \texttt{JMP}, \texttt{JPC}, dan \texttt{CAL} menunjuk ke alamat yang valid. Selama eksekusi, interpreter juga memvalidasi:

\begin{itemize}
    \item Akses stack tidak out-of-bounds untuk \texttt{LOD} dan \texttt{STO}.
    \item Penulisan ke slot frame header dilindungi (\texttt{isProtectedFrameSlot}).
    \item Operasi aritmatika tidak menyebabkan overflow/underflow menggunakan \texttt{checkedAdd}, \texttt{checkedSub}, \texttt{checkedMul}, dan \texttt{checkedNeg}.
    \item Pembagian dan modulo tidak membagi dengan nol.
    \item Kedalaman stack tidak melebihi batas maksimum untuk mencegah stack overflow.
\end{itemize}

\section{Format Output}
Output program Milestone 4 terdiri dari tujuh bagian: \texttt{tab}, \texttt{btab}, \texttt{atab}, parse tree, decorated AST, intermediate code, dan interpreter output (atau interpreter errors jika terjadi runtime error). Intermediate code ditampilkan dalam format baris demi baris dengan nomor instruksi, opcode, level, dan operand. Interpreter output menampilkan hasil eksekusi program, termasuk nilai yang ditulis oleh \texttt{writeln}.
